\section{Conclusion}
\label{sec:concl}

Serverless workflow orchestration has become a backbone of cloud-native systems, yet
its workflows are untyped configuration whose correctness is checked only when they
run. We presented \tool, a static verifier that catches four classes of orchestration
defect ---contract mismatches, out-of-order business transitions, unsafe retries, and
missing compensations--- \emph{before} deployment, and compiles verified workflows
back to executable Amazon States Language. The key to analysing the workflows that
already exist is a two-layer model: sound native checks that need nothing beyond
\asl, paired with a lightweight annotation layer whose defaults are inferred from the
naming and resource conventions that pervade real workflows.

Our evaluation on 193 real Step Functions workflows supports the three research
questions. \tool{} detects the targeted defect classes (100\% recall over 429
injected defects) and surfaces genuine issues in unmodified workflows (RQ1); it
recovers idempotency and persistence at zero authoring cost with 83--91\% accuracy and
visible, correctable residual error (RQ2); and it verifies a workflow in 19\,$\mu$s on
average while adding zero runtime overhead, with verified workflows running unmodified
on real AWS Step Functions (RQ3).

The broader message mirrors the history of type systems: just as static typing moved
classes of error from run time to compile time, lightweight static verification can do
the same for serverless orchestration, on the industrial platforms developers already
use. Several directions remain. Additional frontends (CNCF Serverless Workflow, SAM,
CDK) would widen reach; richer effect and ownership information would sharpen the
contract and compensation analyses; and a precise data-flow analysis of \texttt{Map}
and \texttt{Parallel} item contexts would extend native contract checking. We release
\tool, the corpus, and the evaluation harness to support replication and reuse.
